\section{Detection and Scoring Choices}
\label{sec:detector-scoring}

\subsection{A compact count-based detector}

The detector keeps the four-category DM and Minka estimator. Two complementary
measurements, total packets and ARP packets, come directly from the existing
count vector: total volume is not modelled by a distribution conditioned on
$N$, and a small ARP change can be masked by dominant traffic. No additional
packet fields, model family, dependency, or device-specific count rule is
required.

Fit-only medians and IQRs scale the complete DM log-probability,
$-\log(1+N)$, and $-\log(1+x_{\mathrm{ARP}})$; a constant fit feature uses unit
scale. Their minimum is calibrated once at the configured lower quantile on
separate normal calibration windows. Joint calibration accommodates correlated
signals without label-tuned weights or separate extreme-tail allocations.
Events retain the underlying DM value in \path{composition_score}; their
\path{feature_margins} are signed distances from the calibrated cutoffs. The
combined score is the minimum margin, alerts below zero, and has severity
$\max(0,-s)$. It is not an attack probability.

During adaptation the original DM and count cutoffs remain detection
references. An update cannot suppress an alert supported by the reference DM.
Learning still uses raw DM scores, and conservative admission also requires
safe volume and ARP margins. The shared batch/streaming path and saved-state
contract therefore support the same detector without changing the fitted
category model.

\subsection{Cross-device comparison}

The comparison pools 2,513 labelled windows: Chromecast and Samsung development
and later partitions, plus 475 archived Echo development windows. Each partition
and mode starts independently from the same fitted normal reference. Confusion
counts are summed before calculating pooled F1; this is not a mean of device
F1 values. The retained method uses the same settings across devices and is
compared with DM-only scoring under each learning policy.

\begin{table}[H]
\centering\small
\begin{tabular}{lrr}
\toprule
\textbf{Mode and policy} & \textbf{DM-only F1} & \textbf{Retained detector F1} \\
\midrule
Static & 68.99\% & 70.68\% \\
Threshold, conservative & 64.49\% & 67.94\% \\
Periodic refit, conservative & 68.99\% & 70.68\% \\
Threshold, legacy & 57.01\% & 70.68\% \\
Periodic refit, legacy & 57.67\% & 71.20\% \\
\bottomrule
\end{tabular}
\caption{Retrospective pooled comparisons with a 288-window buffer and
144-window update cadence. These runs are distinct from the historical
288/240 cadence experiment.}
\label{tab:retained-scoring}
\end{table}

Static classifications on Chromecast and Samsung are unchanged. Chromecast
development ROC--AUC improves from 0.9092 to 0.9332; Echo static F1 improves
from 46.15\% to 51.16\%, removing eight false positives while losing one true
positive. Under the legacy learning policy, Samsung development F1 improves
from 54.00\% to 78.69\%, and its later F1 from 66.67\% to 90.91\%.
These are aggregate gains with specific tradeoffs: Echo legacy modes add six
false positives, and the later Chromecast legacy audit retains one original
false alert that adaptation alone suppressed.

All conservative candidates on Chromecast and Samsung remain rejected at their
existing settings. Echo exercises two accepted conservative threshold updates;
its threshold-adaptive F1 improves from 34.29\% to 42.86\%. These results do not
establish a successful conservative model-refitting benefit. D-Link static alert
counts remain 3/270 development and 6/288 later windows; without labels, neither
count is a measured FPR.

\subsection{Rejected choices and scope of the evidence}

Per-packet score normalisation loses attack recall on both reference labelled
devices. Broad upper tails for every category improve some development results
but produce 78 false positives on later Samsung traffic, compared with none
for its DM-only baseline. Separate extreme-tail budgets and a joint rank-tail
cutoff also lose later attack coverage. These choices are excluded. Byte, SYN,
and direction features are not part of the retained contract because their
packet metadata cannot be validated consistently across the labelled datasets.

The evidence supports a small scoring refinement and protection against
adaptation erasing reference evidence, not a universal recall improvement.
A fixed reference can retain false alerts after benign drift, and an empirical
calibration quantile does not guarantee a future false-alert rate. Later labels
have been inspected during candidate assessment, so the reported comparisons
are retrospective rather than an untouched final test. The decision record and
per-device results are indexed in Table~\ref{tab:detection-evidence}.
